\lyseis{7.9}{Δυνάμεις με εκθέτη ακέραιο}
\sol{1} α) $\frac{1}{2}$ \quad β) $\frac{1}{8}$ \quad γ) $\frac{1}{100}$ \quad δ) $\frac{1}{25}$ \quad ε) $1$ \quad στ) $1$.
\sol{2} α) $0{,}1$ \quad β) $0{,}001$ \quad γ) $0{,}00001$.
\sol{3} α) $-\frac{1}{8}$ \quad β) $\frac{1}{4}$ \quad γ) $-1$ \quad δ) $8$ \quad ε) $\frac{9}{4}$ \quad στ) $-\frac{4}{3}$.
\sol{4} α) $3^{-2}$ \quad β) $2^{-5}$ \quad γ) $10^{-4}$ \quad δ) $(-5)^{-3}$.
\sol{5} α) $2^2=4$ \quad β) $10^6$ \quad γ) $\frac{1}{9}$ \quad δ) $\frac{1}{100}$ \quad ε) $7^0=1$ \quad στ) $(-2)^2=4$.
\sol{6} α) $4+1-\frac{1}{2}=\frac{9}{2}$ \quad β) $\frac{1}{4}+\frac{1}{2}=\frac{3}{4}$ \quad γ) $10^{2-5+3}=10^0=1$ \quad δ) $\frac{1}{9}$ \quad ε) $1\cdot\frac{1}{4}=\frac{1}{4}$.
\sol{7} α) $-4$ \quad β) $-3$ \quad γ) $-3$ \quad δ) $-3$ (αφού $\left(\frac{1}{2}\right)^{-3}=8$) \quad ε) $0$.
\sol{8} α) Λάθος: $\frac{1}{2}$. β) Σωστό. γ) Σωστό. δ) Λάθος: $0{,}001$. ε) Σωστό. στ) Σωστό. ζ) Λάθος: $\frac{1}{4}$.
\sol{9} $10^3$ μ., $10^{-2}$ μ., $10^{-3}$ μ., $10^{-6}$ μ., $10^{-9}$ μ.
\sol{10} Μένει $\left(\frac{1}{2}\right)^5=\frac{1}{32}=2^{-5}$.
\sol{11} Μετά από $3$ ημέρες $64\cdot2^{-3}=8$ g. Πριν από $2$ ημέρες $64\cdot2^2=256$ g.
\sol{12} $3^4=81$, $3^3=27$, $3^2=9$, $3^1=3$, $3^0=1$, $3^{-1}=\frac{1}{3}$, $3^{-2}=\frac{1}{9}$, $3^{-3}=\frac{1}{27}$.
\sol{13} $400\cdot2^{-3}=50$ €.
\sol{14} Στις δυνάμεις του $2$ όταν ο εκθέτης μειώνεται κατά $1$ η τιμή γίνεται το μισό: $2^1=2$, $2^0=1$. Το $1$ είναι το μισό του $2$, όχι το $0$. Το ίδιο ισχύει για κάθε βάση.
\sol{15} $\left(\frac{2}{3}\right)^{-2}=\frac{1}{(2/3)^2}=\frac{1}{4/9}=\frac{9}{4}$ και $\left(\frac{3}{2}\right)^2=\frac{9}{4}$.
\sol{16} $2^{-3}=\frac{1}{8}$ (θετικό), ενώ $-2^3=-8$.
\sol{17} $\alpha^{-3}\cdot\alpha^3=\alpha^0$ και, αφού $\alpha^{-3}=\frac{1}{\alpha^3}$, το γινόμενο είναι $1$. Άρα $\alpha^0=1$.
\sol{18} Θα ήταν $\frac{1}{0}$, και δεν διαιρούμε με το $0$.
\sol{19} $-1$, $1$, $-1$. Το $(-1)^{-\nu}=(-1)^\nu$: $1$ για άρτιο και $-1$ για περιττό.
\sol{20} $0{,}1$ mm $=10^{-1}\cdot10^{-3}$ μ. $=10^{-4}$ μ.
